% Generated draft: requires mathematical and source review.
\documentclass[11pt]{article}
\usepackage[margin=25mm]{geometry}
\usepackage{fontspec}
\setmainfont{Latin Modern Roman}
\usepackage{amsmath,amssymb,mathtools,unicode-math}
\setmathfont{Latin Modern Math}
\usepackage{enumitem,needspace,xurl,hyperref}
\setlength{\parindent}{0pt}
\setlength{\parskip}{5pt}
\setlength{\emergencystretch}{3em}
\newcommand{\N}{\mathbb{N}}
\newcommand{\Z}{\mathbb{Z}}
\newcommand{\Q}{\mathbb{Q}}
\newcommand{\R}{\mathbb{R}}
\newcommand{\C}{\mathbb{C}}
\newcommand{\F}{\mathbb{F}}
\newcommand{\A}{\mathbb{A}}
\newcommand{\PP}{\mathbb{P}}
\newcommand{\E}{\mathbb{E}}
\newcommand{\eps}{\varepsilon}
\newcommand{\dd}{\,\mathrm{d}}
\begin{document}
\section*{UnsolvedMath ID: 30005585}
\textbf{Source problem code:} catalogue entry
\subsection{Definitions and mathematical statement}
Write \(\mathbb T^2=\R^2/\Z^2\). A Hamiltonian system consists of a symplectic manifold \((M,\omega)\) and a smooth Hamiltonian \(H:M\to\R\); fix the convention \(\iota_{X_H}\omega=-\dd H\). Its flow can be considered on a specified invariant domain, such as a regular energy hypersurface \(H^{-1}(e)\). The source does not specify such a domain or energy.

For a complete smooth flow \(\phi:\R\times N\to N\), distinguish the conditions
\[
\forall x\in N\;\exists T_x>0:\quad \phi_{T_x}(x)=x
\]
and
\[
\exists T>0\;\forall x\in N:\quad
\phi_T(x)=x\quad\text{and}\quad
\phi_s(x)\ne x\ \text{for every }0<s<T.
\]
The first says that every trajectory is periodic, allowing stationary trajectories; the second says that every trajectory has the same positive least period. The latter is a possible Zoll convention. These conditions must not be identified without further hypotheses. The supplied record does not specify the intended Zoll convention, whether stationary trajectories are excluded, or whether time reparametrization is allowed.

Geodesic flows arise from a Riemannian metric \(g\), usually on a specified positive energy level of the kinetic Hamiltonian \(H_g(q,p)=\tfrac12|p|_g^2\). Magnetic flows additionally involve a closed two-form \(\sigma\) on the configuration manifold, represented in the Hamiltonian description by a twisted cotangent symplectic form. On \(\mathbb T^2\), the metric, magnetic form, and energy are data whose permitted variations must be specified.

\textbf{Source question, with unresolved parameters.} For a specified class \(\mathcal C\) of Hamiltonian flows and a specified Zoll convention, let
\[
\mathcal Z(\mathcal C)=\{\phi\in\mathcal C:\phi\text{ is Zoll under that convention}\}.
\]
Characterize the abundance of \(\mathcal Z(\mathcal C)\) within \(\mathcal C\), with geodesic flows and magnetic flows on the two-torus among the examples mentioned in the supplied clean statement. The source does not determine \(\mathcal C\), its topology, an equivalence relation on systems, or the meaning of abundance, such as density or the existence and dimension of families. Consequently, this is a research question requiring those choices, rather than a uniquely specified quantified conjecture. The original excerpt is truncated or merged with a talk heading, so its intended grouping of the geodesic and two-torus examples also requires review.
\subsection{Short English statement}
The question asks how plentiful Zoll Hamiltonian flows are in specified classes, including geodesic and magnetic examples. A precise formulation requires a definition of Zoll, the allowed systems and energy levels, and a meaning of abundance. Supplied dated curation reports progress for magnetic systems on \(\mathbb T^2\) and labels the broader question partially solved; this draft does not independently verify that assessment.
\subsection{Sources}
\begin{itemize}
\item Local UnsolvedMath registry, permanent ID \(30005585\), source code \texttt{catalogue entry}, title \emph{Abundance of Zoll Flows in Hamiltonian Systems}, Oberwolfach Reports Open Problems collection. The supplied registry labels the statement \texttt{corrected\_verified} and claims comparison with the official report; that verification is supplied provenance, not an independent check here.
\item Supplied report citation: \emph{Dynamische Systeme (2024)}, Oberwolfach Reports, DOI \texttt{10.4171/owr/2023/30}, pp.\ 1671--1730. \url{https://doi.org/10.4171/owr/2023/30}. The record separately gives proposed year 2023; the year discrepancy is unresolved.
\item The original excerpt includes the apparent talk heading \emph{Zoll magnetic systems on the two-torus: a Nash--Moser approach}, naming Gabriele Benedetti, joint work with Luca Asselle and Massimiliano Berti. The registry supplies no explicit proposer; the heading alone does not establish who proposed the broader question.
\item Supplied literature reference: Asselle--Benedetti--Berti, \url{https://doi.org/10.1016/j.aim.2024.109826}.
\item Supplied literature reference: Asselle--Benedetti, \url{https://arxiv.org/abs/1909.13821}.
\item Literature assessment dated exactly \texttt{2026-08-22}: supplied curation labels the problem partially solved and reports an infinite-dimensional family of Zoll magnetic systems on the two-torus near flat or constant-field examples, substantially answering abundance in that class, while the broader Hamiltonian and geodesic-flow question remains open. This is supplied dated curation derived from the OWR record, not independently verified literature assessment. The precise hypotheses and scope of the cited results require checking.
\end{itemize}
\end{document}
